\section{Introduction}
\label{sec:intro}
This project creates a complete, time-consistent 3D video of a traffic intersection from eight synchronized cameras to audit the line-of-sight between pedestrians and approaching vehicles \cite{mildenhall2020nerf}. A two-pass method first generates a 3D snapshot for each moment, then temporally fuses partial views of dynamic subjects into complete 3D models so that occluded geometry is recovered before safety analysis, extending multi-view reconstruction ideas popularized by neural radiance fields~\cite{mildenhall2020nerf}. The final interactive 3D scene allows for precise ray-casting analysis to determine when a pedestrian's view of any vehicle travelling toward the intersection was obstructed, closing the gap between manual sight-distance checks and data-driven audits~\cite{aashto2018greenbook}.

\subsection{Problem Statement}
At urban intersections, pedestrians at the curb often face blocked views of approaching vehicles due to occlusions like parked cars, poles, or other obstacles. This visibility issue creates a serious safety risk, impairing people's ability to judge traffic conditions before crossing and contributing to preventable accidents documented in intersection safety reviews~\cite{lord2007safety}. In the United States alone, 7,388 pedestrians were killed in traffic crashes in 2021, with limited sight distance called out as a contributing factor in urban corridors~\cite{nhtsa2023pedestrian}.

\subsection{Research Significance}
Solving this matters for safer streets and better urban design. By building time-consistent 3D models from multi-camera data, this project audits line-of-sight (LoS), a core principle of transportation safety \cite{aashto2018policy}, between pedestrians and vehicles, pinpointing exact blockers that hide dangers. These insights can drive real fixes, such as removing or redesigning obstacles, adding signs, or implementing smart traffic controls. Broader impacts include fewer pedestrian-vehicle crashes, improved transportation planning, and more accessible cities. Quantitatively, enhanced LoS could reduce accident rates, cut costs in healthcare and insurance, and minimize disruptions in daily mobility, leading to more resilient and people-centered urban environments while keeping documentation compatible with transportation design guidance~\cite{aashto2018greenbook}.
